% !PRINTED 503
\weberpage{503}

Da, wie wir oben schon gesehen haben, zwei Nebengruppen zu $Q$ entweder ganz identisch sind, oder kein gemeinschaftliches Element enthalten, so können wir noch den Satz aussprechen:

2. \gesperrt{Ist $\pi$ eine beliebige Permutation aus $P$, so ist das System der Nebengruppen}
\[
Q\pi,\ Q\pi_1\pi,\ Q\pi_2\pi\ldots Q\pi_{j-1}\pi
\]
\gesperrt{von dem System}
\[
Q,\ Q\pi_1,\ Q\pi_2\ldots Q\pi_{j-1}
\]
\gesperrt{nur durch die Anordnung unterschieden.}

Wir dehnen die Bezeichnung Nebengruppen auch auf Systeme von der Form $\pi Q$ aus, und können $P$ auch in solche Systeme zerlegen. Es folgt speciell aus der Zerlegung (4) die zweite
\[
P = Q + \pi_1^{-1}Q + \pi_2^{-1}Q + \cdots \pi_{j-1}^{-1}Q.
\]
Denn wenn $\pi_1$ nicht in $Q$ vorkommt, so kommt auch $\pi_1^{-1}$ nicht darin vor, und wir erhalten eine erste Nebengruppe $\pi_1^{-1}Q$. Wenn dann $\pi_2$ nicht in $Q\pi_1$ vorkommt, so kommt auch $\pi_2^{-1}$ nicht in $\pi_1^{-1}Q$ vor; denn aus $\pi_2^{-1} = \pi_1^{-1}\varkappa$ würde $\pi_2 = \varkappa^{-1}\pi_1$ folgen u.\,s.\,f.

Wir geben hier der Zerlegung (4) den Vorzug wegen der Bedeutung, die sie, wie wir sogleich sehen werden, für das algebraische Problem hat.

Wir nennen eine Grösse $\psi$ des Körpers $N$ \gesperrt{zu der Gruppe $Q$ gehörig} (wenn nöthig mit dem Zusatz \glqq{}innerhalb $P$\grqq{}), wenn $\psi$ sich nicht ändert, falls eine der Operationen von $Q$ darauf angewandt wird, dagegen sich ändert, wenn irgend eine andere Operation aus $P$ angewandt wird. Wie man sieht, ist dieser Begriff der Zugehörigkeit gegen den im \S.\,152 festgesetzten etwas erweitert, insofern die Permutationen, die ausserhalb $P$ liegen, hier gar nicht in Betracht kommen. Für den Fall, dass $P$ die symmetrische Gruppe ist, fallen aber beide Begriffsbestimmungen zusammen.

Bei dieser Definition der Zugehörigkeit gelten nun die folgenden Sätze allgemein.

3. \gesperrt{Zu jedem Theiler $Q$ von $P$ gehören Grössen in $N$, und jede Grösse in $N$ gehört zu einem bestimmten Theiler $Q$ von $P$.}

Der erste Theil des Satzes ist bereits im \S.\,152, 5. bewiesen. Denn wenn eine Function der Grössen $\alpha$ innerhalb der symme-
%CONT>>
